% ============================================================
%  00_CARATULA_DOCX.TEX
%  Versión de la carátula SOLO para el flujo a Word
%  (convertir_docx.sh la usa en vez de secciones/00_caratula.tex).
%  ------------------------------------------------------------
%  Por qué existe un archivo aparte: secciones/00_caratula.tex usa
%  maquetación de bajo nivel (\fontsize{20}{26}, tres \minipage en
%  fila) que Pandoc no traduce a Word --- se comprobó que Pandoc
%  ni siquiera respeta \begin{center} o \Large/\LARGE en un .docx
%  (los descarta en silencio). La fila de logos SÍ se resuelve bien
%  como tabla de 3 columnas sin bordes (Pandoc convierte tabular en
%  tabla real de Word). El tamaño y centrado del resto del texto
%  NO se resuelve aquí: lo aplica docx/estilizar_caratula.py después
%  de la conversión, buscando cada párrafo por su posición --- ver
%  ese script para la lista exacta de párrafos que ajusta.
%
%  Si cambias el CONTENIDO de esta carátula, mantén un párrafo por
%  bloque (sin líneas en blanco de más) para que las posiciones que
%  usa estilizar_caratula.py sigan siendo válidas.
% ============================================================
\noindent
\begin{tabular}{p{0.25\textwidth}p{0.5\textwidth}p{0.25\textwidth}}
\includegraphics[height=1.3cm]{imagenes/logo_umss.png}\includegraphics[height=1.3cm]{imagenes/logo_facultad.png}
&
UNIVERSIDAD MAYOR DE SAN SIMÓN \newline FACULTAD DE CIENCIAS Y TECNOLOGÍA \newline DIRECCIÓN DE POSGRADO
&
\includegraphics[height=1.3cm]{imagenes/logo_posgrado.png}
\end{tabular}

\textbf{\MakeUppercase{\tituloMonografia}}

\textbf{MONOGRAFÍA PRESENTADA PARA OBTENER EL CERTIFICADO DE DIPLOMADO EN \MakeUppercase{\nombreDiplomado}\ \nombreVersion\ VERSIÓN}

\textbf{POSTULANTE: \MakeUppercase{\nombrePostulante}}

\textbf{Cochabamba -- Bolivia}

\textbf{\anioMonografia}

\clearpage
